\subsubsection*{Question a) Distance par rapport à la sirène de but}

\begin{enumerate}[label=\arabic*)]
    \item \textbf{Variables connues et recherchées :}
    \begin{itemize}[label= -]
        \item Puissance de la source : $P = \qty{\v{PneuPuissance}}{W}$
        \item Niveau d'intensité sonore : $\beta_{\text{sirène}} = \qty{\v{PneuNiveaudB}}{dB}$
        \item Intensité du seuil d'audibilité : $I_0 = 10^{-12} \text{ W/m}^2$
        \item Variable recherchée : Distance $r \text{ (m)}$
    \end{itemize}

    \item \textbf{Équations pertinentes :}
    \[
    \beta = 10 \log_{10}\left(\frac{I}{I_0}\right) \quad \implies \quad I = I_0 \cdot 10^{\frac{\beta}{10}}
    \]
    \[
    I = \frac{P}{4\pi r^2} \quad \text{(modèle d'onde sphérique)}
    \]

    \item \textbf{Isolation de la variable recherchée :}
    \[
    \frac{P}{4\pi r^2} = I_0 \cdot 10^{\frac{\beta}{10}} \quad \implies \quad r^2 = \frac{P}{4\pi \cdot I_0 \cdot 10^{\frac{\beta}{10}}} \quad \implies \quad r = \sqrt{\frac{P}{4\pi \cdot I_0 \cdot 10^{\frac{\beta}{10}}}}
    \]

    \item \textbf{Résolution numérique :}
    \[
    I_{\text{sirène}} = 10^{-12} \cdot 10^{\frac{\v{PneuNiveaudB}}{10}} = \num[scientific-notation=true, round-mode=figures, round-precision=3]{\fpeval{10^(-12) * 10^(\v{PneuNiveaudB} / 10)}} \text{ W/m}^2
    \]
    \[
    r = \sqrt{\frac{\v{PneuPuissance}}{4\pi \cdot 10^{-12} \cdot 10^{\frac{\v{PneuNiveaudB}}{10}}}} = \qty{\fpeval{round(sqrt(\v{PneuPuissance} / (4 * pi * 10^(-12) * 10^(\v{PneuNiveaudB} / 10))), 2)}}{m}
    \]

    \item \textbf{Réponse :}
    Vous vous trouvez à une distance d'environ \textbf{\qty{\fpeval{round(sqrt(\v{PneuPuissance} / (4 * pi * 10^(-12) * 10^(\v{PneuNiveaudB} / 10))), 2)}}{m}} de la sirène de but.
\end{enumerate}


% ------------------------------------------------------------------------------
% QUESTION B
% ------------------------------------------------------------------------------
\subsubsection*{Question b) Puissance acoustique absorbée par le tympan}

\begin{enumerate}[label=\arabic*)]
    \item \textbf{Variables connues et recherchées :}
    \begin{itemize}[label= -]
        \item Rayon du conduit auditif : $R = \qty{\v{OreilleRayon}}{m}$
        \item Intensité sonore locale : $I_{\text{sirène}} = \num[scientific-notation=true, round-mode=figures, round-precision=3]{\fpeval{10^(-12) * 10^(\v{PneuNiveaudB} / 10)}} \text{ W/m}^2$
        \item Variable recherchée : Puissance absorbée $P_{\text{tympan}} \text{ (W)}$
    \end{itemize}

    \item \textbf{Équations pertinentes :}
    \[
    A = \pi R^2 \quad \text{(surface circulaire)}
    \]
    \[
    P_{\text{tympan}} = I_{\text{sirène}} \cdot A
    \]

    \item \textbf{Isolation de la variable recherchée :}
    \[
    P_{\text{tympan}} = \left( I_0 \cdot 10^{\frac{\beta}{10}} \right) \cdot \pi R^2
    \]

    \item \textbf{Résolution numérique :}
    \[
    A = \pi \cdot (\v{OreilleRayon})^2 = \num[scientific-notation=true, round-mode=figures, round-precision=3]{\fpeval{pi * (\v{OreilleRayon})^2}} \text{ m}^2
    \]
    \[
    P_{\text{tympan}} = \left(10^{-12} \cdot 10^{\frac{\v{PneuNiveaudB}}{10}}\right) \cdot \pi (\v{OreilleRayon})^2 = \qty[scientific-notation=true, round-mode=figures, round-precision=3]{\fpeval{(10^(-12) * 10^(\v{PneuNiveaudB} / 10)) * pi * (\v{OreilleRayon})^2}}{W}
    \]

    \item \textbf{Réponse :}
    La puissance acoustique exacte absorbée par le tympan est de \textbf{\qty[scientific-notation=true, round-mode=figures, round-precision=3]{\fpeval{(10^(-12) * 10^(\v{PneuNiveaudB} / 10)) * pi * (\v{OreilleRayon})^2}}{W}} (soit environ \qty{\fpeval{round((10^(-12) * 10^(\v{PneuNiveaudB} / 10)) * pi * (\v{OreilleRayon})^2 * 10^6, 2)}}{\micro\watt}).
\end{enumerate}


% ------------------------------------------------------------------------------
% QUESTION C
% ------------------------------------------------------------------------------
\subsubsection*{Question c) Niveau d'intensité sonore total (Superposition)}

\begin{enumerate}[label=\arabic*)]
    \item \textbf{Variables connues et recherchées :}
    \begin{itemize}[label=-]
        \item Niveau sonore de la sirène : $\beta_{\text{sirène}} = \qty{\v{PneuNiveaudB}}{dB}$
        \item Niveau sonore de la musique : $\beta_{\text{musique}} = \qty{\v{MusiqueNiveaudB}}{dB}$
        \item Intensité de référence : $I_0 = 10^{-12} \text{ W/m}^2$
        \item Variable recherchée : Niveau sonore total $\beta_{\text{tot}} \text{ (dB)}$
    \end{itemize}

    \item \textbf{Équations pertinentes :}
    \[
    I_{\text{tot}} = I_{\text{sirène}} + I_{\text{musique}} = I_0 \cdot 10^{\frac{\beta_{\text{sirène}}}{10}} + I_0 \cdot 10^{\frac{\beta_{\text{musique}}}{10}}
    \]
    \[
    \beta_{\text{tot}} = 10 \log_{10}\left(\frac{I_{\text{tot}}}{I_0}\right)
    \]

    \item \textbf{Isolation de la variable recherchée :}
    \[
    \beta_{\text{tot}} = 10 \log_{10}\left(\frac{I_0 \left(10^{\frac{\beta_{\text{sirène}}}{10}} + 10^{\frac{\beta_{\text{musique}}}{10}}\right)}{I_0}\right) = 10 \log_{10}\left(10^{\frac{\beta_{\text{sirène}}}{10}} + 10^{\frac{\beta_{\text{musique}}}{10}}\right)
    \]

    \item \textbf{Résolution numérique :}
    \[
    I_{\text{musique}} = 10^{-12} \cdot 10^{\frac{\v{MusiqueNiveaudB}}{10}} = \num[scientific-notation=true, round-mode=figures, round-precision=3]{\fpeval{10^(-12) * 10^(\v{MusiqueNiveaudB} / 10)}} \text{ W/m}^2
    \]
    \[
    I_{\text{tot}} = \num[scientific-notation=true, round-mode=figures, round-precision=3]{\fpeval{10^(-12) * (10^(\v{PneuNiveaudB} / 10) + 10^(\v{MusiqueNiveaudB} / 10))}} \text{ W/m}^2
    \]
    \[
    \beta_{\text{tot}} = 10 \log_{10}\left(10^{\frac{\v{PneuNiveaudB}}{10}} + 10^{\frac{\v{MusiqueNiveaudB}}{10}}\right) = \qty{\fpeval{round(10 * (ln(10^(\v{PneuNiveaudB} / 10) + 10^(\v{MusiqueNiveaudB} / 10)) / ln(10)), 2)}}{dB}
    \]

    \item \textbf{Réponse :}
    Le niveau d'intensité sonore total ressenti lorsque la sirène retentit en même temps que la musique est de \textbf{\qty{\fpeval{round(10 * (ln(10^(\v{PneuNiveaudB} / 10) + 10^(\v{MusiqueNiveaudB} / 10)) / ln(10)), 2)}}{dB}}.
\end{enumerate}